% Paper Tables: New experimental results
% Confusion Matrix, Reachability Separation, Compiler Comparison,
% Cross-Constellation, Polar Exclusion, Runtime Benchmark

% ============================================================
% Table: Verifier Confusion Matrix (3-way)
% ============================================================
\begin{table}[t]
\centering
\caption{Three-way validator confusion matrix on the 240-intent benchmark
(8-pass pipeline with constructive feasibility certification).
\textsc{Accept} requires a routing witness; \textsc{Abstain} indicates
unsupported constraint combinations (safe non-accept). The validator
achieves 0\% unsafe acceptance across all categories.}
\label{tab:confusion_matrix}
\begin{tabular}{lccc}
\toprule
& \textbf{Accept} & \textbf{Reject} & \textbf{Abstain} \\
\midrule
Single ($n$=80)         & 10  & 5   & 65  \\
Compositional ($n$=100) & 40  & 25  & 35  \\
Conditional ($n$=30)    & 8   & 2   & 20  \\
Infeasible ($n$=30)     & 0   & 22  & 8   \\
\midrule
\textbf{Total} ($n$=240) & 58 & 54 & 128 \\
\midrule
\multicolumn{4}{l}{\textit{Safety (unsafe = infeasible \textsc{Accept}):}} \\
\quad Infeasible       & \multicolumn{3}{l}{0/30 = 0\% (was 72\% with 7-pass)} \\
\multicolumn{4}{l}{\textit{Coverage (\textsc{Accept}+\textsc{Reject}):}} \\
\quad Decided          & \multicolumn{3}{l}{112/240 = 46.7\%} \\
\bottomrule
\multicolumn{4}{l}{\footnotesize 32 feasible programs rejected as routing-infeasible;} \\
\multicolumn{4}{l}{\footnotesize 29/32 independently confirmed via separate Dijkstra oracle.} \\
\end{tabular}
\end{table}

% ============================================================
% Table: Reachability Separation
% ============================================================
\begin{table}[t]
\centering
\caption{Reachability separation analysis. Raw PDR differences in
polar/compositional scenarios are entirely explained by the reachability
ceiling (76\% of OD pairs unreachable), not routing quality. Both GNN
and Dijkstra achieve 100\% delivery on reachable pairs.}
\label{tab:reachability}
\begin{tabular}{lccccc}
\toprule
\textbf{Scenario} & \textbf{Reach.} & \multicolumn{2}{c}{\textbf{Raw PDR}} & \multicolumn{2}{c}{\textbf{Reachable PDR}} \\
\cmidrule(lr){3-4} \cmidrule(lr){5-6}
& & GNN & Dijkstra & GNN & Dijkstra \\
\midrule
Baseline          & 100\%  & 99.8\% & ---     & 99.8\% & ---     \\
Node failure      & 100\%  & 98.7\% & 97.8\%  & 98.7\% & 97.8\% \\
Plane maint.      & 100\%  & 70.5\% & 70.2\%  & 70.5\% & 70.2\% \\
Polar avoid.      & 24.0\% & 34.6\% & 47.9\%  & 100\%  & 100\%  \\
Compositional     & 24.0\% & 34.3\% & 47.1\%  & 100\%  & 100\%  \\
\bottomrule
\end{tabular}
\end{table}

% ============================================================
% Table: Compiler Comparison (Rule-Based vs LLM)
% ============================================================
\begin{table}[t]
\centering
\caption{Intent compiler comparison on the 240-intent benchmark
(193 feasible intents; 17 additional routing-infeasible discovered
by Pass~8 excluded from feasible metrics).
The rule-based parser handles simple intents adequately but degrades
sharply on compositional intents, confirming that intent compilation
is a compositional reasoning task that benefits from LLM capabilities.}
\label{tab:compiler_comparison}
\begin{tabular}{lccc}
\toprule
\textbf{Metric} & \textbf{Rule-Based} & \textbf{LLM 4B} & \textbf{LLM 9B} \\
\midrule
Compiled          & 100.0\% & 59.6\% & \textbf{98.4\%} \\
Types Match       & 67.1\%  & 55.4\% & \textbf{91.7\%} \\
Full Match        & 56.7\%  & 54.2\% & \textbf{87.6\%} \\
Avg Latency       & \textbf{0.05ms} & 204s & 15.7s \\
\midrule
\multicolumn{4}{l}{\textit{Full match by category:}} \\
\quad Single      & 76.2\%  & ---    & \textbf{89.5\%} \\
\quad Compositional & 40.0\% & ---   & \textbf{86.2\%} \\
\quad Conditional & 66.7\%  & ---    & \textbf{86.7\%} \\
\quad Infeasible  & 50.0\%  & ---    & \textbf{73.3\%} \\
\bottomrule
\end{tabular}
\end{table}

% ============================================================
% Table: OOD Generalization
% ============================================================
\begin{table}[t]
\centering
\caption{Out-of-distribution generalization on paraphrased intents
(38 total: 33 scorable + 5 ambiguous). The compiler maintains 81.8\%
full match accuracy with only 4.4pp degradation from template intents,
demonstrating robust generalization to novel phrasings.}
\label{tab:ood}
\begin{tabular}{lccc}
\toprule
\textbf{Category} & \textbf{N} & \textbf{Compiled} & \textbf{Full Match} \\
\midrule
Single         & 20 & 100\% & 95.0\% (19/20) \\
Compositional  & 5  & 100\% & 40.0\% (2/5)   \\
Conditional    & 8  & 100\% & 75.0\% (6/8)   \\
Ambiguous      & 5  & 100\% & qualitative: 5/5 \\
\midrule
Scorable total & 33 & 100\% & \textbf{81.8\%} (27/33) \\
\bottomrule
\end{tabular}
\end{table}

% ============================================================
% Table: Cross-Model Scaling (4B vs 9B)
% ============================================================
\begin{table}[t]
\centering
\caption{Cross-model scaling: 4B vs 9B parameter LLM on the full
240-intent benchmark. The 9B model dramatically outperforms the 4B
model, suggesting a compositional complexity threshold between 4B
and 9B parameters for this task.}
\label{tab:cross_model}
\begin{tabular}{lcc}
\toprule
\textbf{Metric} & \textbf{Qwen 4B} & \textbf{Qwen 9B} \\
\midrule
Compiled       & 59.6\% & \textbf{98.4\%} \\
Types Match    & 55.4\% & \textbf{91.7\%} \\
Full Match     & 54.2\% & \textbf{87.6\%} \\
First-try Rate & 47.1\% & \textbf{89.2\%} \\
Avg Latency    & 204.4s & \textbf{15.7s}  \\
\bottomrule
\end{tabular}
\end{table}

% ============================================================
% Table: Topology Sweep (GNN vs Dijkstra across severity)
% ============================================================
\begin{table}[t]
\centering
\caption{GNN vs Dijkstra PDR across topology degradation severity
(orbital plane removal, 20 sats/plane). The GNN router matches
Dijkstra within 0.22pp at all severity levels, confirming robust
distillation quality. Non-monotonic PDR reflects topology-dependent
reachability under random plane selection (3-seed average).}
\label{tab:topology_sweep}
\begin{tabular}{rrccc}
\toprule
\textbf{Planes Off} & \textbf{Capacity} & \textbf{GNN PDR} & \textbf{Dijkstra PDR} & \textbf{$\Delta$} \\
\midrule
1  & 5\%  & 81.1\% & 80.9\% & +0.22 \\
2  & 10\% & 41.1\% & 41.1\% & 0.00 \\
3  & 15\% & 36.8\% & 36.8\% & 0.00 \\
5  & 25\% & 45.3\% & 45.3\% & 0.00 \\
7  & 35\% & 42.5\% & 42.5\% & 0.00 \\
10 & 50\% & 11.4\% & 11.4\% & 0.00 \\
13 & 65\% & 5.4\%  & 5.4\%  & 0.00 \\
15 & 75\% & 5.5\%  & 5.5\%  & 0.00 \\
17 & 85\% & 36.4\% & 36.4\% & 0.00 \\
\bottomrule
\end{tabular}
\end{table}

% ============================================================
% Table: Cross-Constellation GNN Generalization
% ============================================================
\begin{table}[t]
\centering
\caption{Zero-shot cross-constellation GNN generalization. The GNN
(trained on Walker Delta 550\,km/53$^\circ$) generalizes perfectly to
altitude changes (same grid topology, different edge weights) but
collapses on SSO 97$^\circ$ where near-polar inclination fundamentally
alters ISL geometry. This motivates the compile--verify--ground pipeline
as the constellation-agnostic safety layer, with the GNN as a fast
accelerator for a specific configuration.}
\label{tab:cross_constellation}
\begin{tabular}{lccc}
\toprule
\textbf{Configuration} & \textbf{GNN PDR} & \textbf{Dijkstra PDR} & \textbf{$\Delta$} \\
\midrule
550\,km / 53$^\circ$ (training) & 99.75\% & 99.75\% & 0.00 \\
1200\,km / 53$^\circ$ (OOD)     & 99.75\% & 99.75\% & 0.00 \\
550\,km / 97$^\circ$ SSO (OOD)  & 45.18\% & 99.09\% & $-$53.91 \\
\bottomrule
\end{tabular}
\end{table}

% ============================================================
% Table: Polar Exclusion Zone GNN Robustness
% ============================================================
\begin{table}[t]
\centering
\caption{GNN robustness under polar exclusion zones. Inter-plane ISLs
are disabled when either endpoint exceeds the latitude threshold.
The GNN (trained with threshold 75$^\circ$, i.e., no exclusion at
53$^\circ$ inclination) degrades proportionally to edge removal,
while Dijkstra maintains 99.75\% PDR throughout. This confirms the
GNN learns topology-specific shortcuts rather than general routing
principles, reinforcing the need for Dijkstra fallback under
topology perturbation.}
\label{tab:polar_exclusion}
\begin{tabular}{rrccc}
\toprule
\textbf{Threshold} & \textbf{Edges Removed} & \textbf{GNN PDR} & \textbf{Dijkstra PDR} & \textbf{$\Delta$} \\
\midrule
30$^\circ$ & 28.8\% & 38.17\% & 99.75\% & $-$61.58 \\
40$^\circ$ & 20.5\% & 54.08\% & 99.75\% & $-$45.67 \\
45$^\circ$ & 15.8\% & 65.20\% & 99.75\% & $-$34.55 \\
50$^\circ$ & 9.8\%  & 80.37\% & 99.75\% & $-$19.38 \\
\bottomrule
\end{tabular}
\end{table}

% ============================================================
% Table: Pass 8 Runtime Benchmark
% ============================================================
\begin{table}[t]
\centering
\caption{8-pass validator runtime on the 240-intent benchmark.
Median verification completes in under 1\,ms; even the worst case
(compositional programs with flow selectors requiring Dijkstra/max-flow
witnesses) stays below 2\,ms. This confirms negligible overhead for
real-time deployment alongside the GNN router (17$\times$ Dijkstra
speedup at inference).}
\label{tab:runtime}
\begin{tabular}{lcccc}
\toprule
\textbf{Category} & \textbf{$n$} & \textbf{Median} & \textbf{P95} & \textbf{Max} \\
\midrule
All programs        & 240 & 0.720\,ms & 1.580\,ms & 1.898\,ms \\
With flow selectors & 98  & 1.059\,ms & 1.738\,ms & 1.898\,ms \\
Topology-only       & 142 & 0.437\,ms & 1.501\,ms & 1.629\,ms \\
\midrule
\multicolumn{5}{l}{\textit{By certification status:}} \\
\quad Accepted      & 58  & 1.044\,ms & 1.738\,ms & 1.812\,ms \\
\quad Rejected      & 32  & 1.135\,ms & 1.757\,ms & 1.898\,ms \\
\quad Abstain       & 128 & 0.428\,ms & 1.507\,ms & 1.629\,ms \\
\bottomrule
\end{tabular}
\end{table}
